\section{Síntesis y ampliación}

\subsection{La cadena completa de Behavior a Mechanism}

Lo más reutilizable de esta clase no es un attention heatmap en particular, sino un procedimiento de investigación: primero se define una métrica de comportamiento y después se busca un abrupt change; se elige un toy model resoluble para construir el álgebra; a partir de la path structure se propone un algorithm candidato; por último, se gradúa la evidencia con ablation, timing, cross-example generality y contraejemplos. La credibilidad de la Mechanistic interpretability proviene de esta cadena, no de una figura que «parece estar atendiendo al token correcto».

\begin{importantbox}{Quince conclusiones centrales}
\begin{enumerate}
  \item La Mechanistic interpretability intenta descompilar los parameters en algorithms.
  \item ICL fija los parámetros; la adaptación temporal ocurre solo en el context y las activations.
  \item La loss surface bidimensional muestra a la vez el training snapshot y la token position.
  \item Cuanto más negativo es $\mathrm{Loss@500}-\mathrm{Loss@50}$, mayor es la late-context advantage.
  \item $0.4$ nats equivalen aproximadamente a $0.58$ bits/token, no a 40\% de accuracy.
  \item La loss bump y el derivative spike dan lugar a la phase-change hypothesis, pero no la demuestran.
  \item Los attention-only toy models eliminan MLP/LayerNorm/bias a cambio de un álgebra exacta.
  \item Una head puede dividirse en dos circuits: QK routing y OV writing.
  \item Una vez fijado el attention pattern, el attention-only path es localmente lineal respecto a la entrada.
  \item Los one-layer models implementan principalmente reglas skip-trigram/positional; su ICL es débil y no presentan phase change.
  \item El eigenvalue lens puede resumir small token-to-token maps, pero deja de funcionar con MLP y en modelos grandes.
  \item La two-layer expansion produce composed/virtual heads y permite que el routing lea features de capas anteriores.
  \item El induction circuit implementa un context nearest neighbor de la forma $[A][B]\ldots[A]\to[B]$.
  \item La small-model ablation es evidencia causal; el formation timing en large models es principalmente evidencia correlacional.
  \item Induction es un mecanismo importante, pero no necesariamente el único; la soft induction y los distributed circuits aún requieren una verificación más sólida.
\end{enumerate}
\end{importantbox}

\subsection{Un Circuit Audit Protocol reutilizable}

Ante una nueva afirmación de interpretabilidad, conviene revisar en orden: si el behavior objetivo está cuantificado; si el component candidato tiene la capacidad estructural necesaria; si el camino se compone a través de layers; si el attention pattern y el OV logit effect son consistentes; si la ablation conserva la distribution y los demás patterns; si el resultado se sostiene entre prompts/seeds/models; y si la conclusión distingue entre toy model, small model y production-scale model.

\begin{knowledgebox}{Lista de verificación práctica}
\begin{enumerate}
  \item Definir la behavior metric y la sign convention;
  \item trazar el diagnóstico bidimensional training × context, en lugar de mirar solo el checkpoint final;
  \item desarrollar los residual paths y distinguir entre routing y writing;
  \item usar activation/attention para formular la hypothesis y ablation para la verificación causal;
  \item registrar compensation, redundancy y negative-result heads;
  \item en la large-model extrapolation, indicar si se trata de correlation, plausibility o intervention;
  \item buscar activamente alternative mechanisms y failure cases.
\end{enumerate}
\end{knowledgebox}

\subsection{Lecturas adicionales}

Se recomienda leer primero la one-/two-layer decomposition de \emph{A Mathematical Framework for Transformer Circuits} y después la evidence taxonomy y el appendix de \emph{In-context Learning and Induction Heads}. A continuación, contrastarlos con las Kaplan scaling laws y el few-shot behavior de GPT-3 para reflexionar sobre la relación entre la behavioral curve, las training dynamics y la circuit formation. Al leer trabajos actuales de interpretación de modelos grandes, conviene observar sobre todo cómo abordan las MLP nonlinearities, la feature superposition, la normalization, la distributed computation y la causal validation, en lugar de comparar únicamente qué tan atractivas son las head visualizations.

\teachervoice{La present theory con la que cierra esta clase consiste en dos algorithms principales: copying generalizado e induction-head in-context nearest neighbors, reconociendo al mismo tiempo que puede haber otros mecanismos. Conservar esta incertidumbre es precisamente la síntesis más fiel a la clase y a la investigación original.}

\end{document}
